\documentclass[10pt,a4paper]{amsart}
\usepackage[margin=19mm]{geometry}
\usepackage{amsmath,amssymb,mathtools}
\usepackage[scr=rsfs,cal=euler]{mathalfa}
\usepackage{xfrac}
\usepackage[unicode,hidelinks,bookmarks=false]{hyperref}
\newcommand{\ie}{i.e.\ }
\newcommand{\cf}{cf.\ }
\newcommand{\resp}{respectively\ }
\newcommand{\Subsection}{\subsection}
\providecommand{\nobreakdash}{\nobreak\hbox{-}}
\setlength{\emergencystretch}{2em}
\multlinegap0pt
\usepackage{xcolor}
\usepackage{enumerate}
\usepackage{bm}
\newtheorem{prop}{Proposition}
\DeclareMathOperator{\M}{\mathbb{M}}
\DeclareMathOperator{\W}{\mathbb{W}}
\DeclareMathOperator{\inv}{\mathrm{inv}}
\title{\bf Algebraic $\bm{n}$-Valued Monoids on $\bm{\mathbb{C}P^1}$,\\ Discriminants and Projective Duality}
\author{Victor Buchstaber, Mikhail Kornev}
\date{Source: arXiv:2510.14010v2 (CC BY 4.0)}
\begin{document}
\maketitle

\noindent\emph{Source and licence.} \bf Algebraic $\bm{n}$-Valued Monoids on $\bm{\mathbb{C}P^1}$,\\ Discriminants and Projective Duality. Version arXiv:2510.14010v2 (CC BY 4.0), distributed by arXiv under the Creative Commons Attribution 4.0 International licence, http://creativecommons.org/licenses/by/4.0/, as recorded in the arXiv licence metadata for this exact version and shown on the abstract page https://arxiv.org/abs/2510.14010v2. Full licence notice: \texttt{licenses/arxiv-2510.14010-CC-BY-4.0.txt}.\par\smallskip

\noindent\textit{Editorial scope.} Counted unit: the article's prop doubling\_monoid (source0.tex lines 781--783). The article's complete original proof of it is delivered with it (source0.tex lines 785--814, one \texttt{proof} environment of 30 lines). No mathematics is written, repaired or re-derived: the statement and the proof are the article's own lines, in the article's own notation, and no retained line is substituted or re-flowed.  No further source-local context is needed: every cross-reference of every retained line resolves inside the two delivered spans.  Nothing was omitted from any retained span, so the package carries no omission record.  No retained line contains a citation, so the wrapper carries no bibliography and no bibliography entry.  No retained line contains a figure environment, an image-inclusion command or drawing code, so no image is delivered and the package declares no asset dependency.  Editorial formatting only, in addition: an AMS \texttt{amsart} preamble that restates the article's own declarations and macros used by the retained lines, namely the environments \texttt{prop} on separate counters, together with the standard packages those lines need.  The retained environments print in this document as Proposition 1 in order of appearance, whereas the article prints the same items with the numbering of its own full document.  The complete native source of the article as distributed in the arXiv e-print for this exact version is bundled unchanged as \texttt{sources/187151/source0.tex}.
\begin{prop}\label{doubling_monoid}
The set $\W$ of doubling elements of a two-valued monoid $($respectively, group$)$ $\M$ forms a diagonal two-valued submonoid $($respectively, subgroup$)$.
\end{prop}

\begin{proof}
Let $\mathbb{W}$ denote the subset of $\M$ consisting of doubling elements. Let $w_1, w_2\in\mathbb{W}$ and let $w_1\ast w_2 = [w_3, w_3]$ for some $w_3\in\M$. Let $m\in\M$ be an arbitrary element. We have
$$\begin{aligned}
[m\ast w_3, m\ast w_3] &= m\ast (w_1\ast w_2)\\
&= (m\ast w_1)\ast w_2.
\end{aligned}$$
Hence $w_3\in \mathbb{W}$, since the multiset $(m\ast w_1)\ast w_2$ is some quadruple point. The case of $w_3\ast m$ is considered similarly. Consider on the set $\mathbb{W}$ the operation
\begin{equation}
\begin{aligned}\label{operation_doubling_submonoids}
\mathbb{W}\times \mathbb{W}&\to \mathbb{W}\\
w_1\cdot w_2 &= w_3.
\end{aligned}
\end{equation}
It is easy to see that the two-valued submonoid $\mathbb{W}\subset\M$ is the two-diagonal of a single-valued monoid $\W$ with operation (\ref{operation_doubling_submonoids}).

{Now let $\M$ be a two-valued group and let $x\in \W$ be a doubling element. We show that $\inv(x)$ is also doubling. For every $y\in\M$:
\begin{equation}\label{invx_is_iterating_element}
x\ast(\inv(x)\ast y) = (x\ast \inv(x))\ast y = [e,e]\ast y =  [y,y,y,y].
\end{equation}
Let $\inv(x)\ast y = [a, b]$ for some $a, b\in \M$. Then from \eqref{invx_is_iterating_element} it follows that
\begin{equation}\label{ax_equals_xb}
x\ast a = x\ast b = [y,y].
\end{equation}
Multiplying equality \eqref{ax_equals_xb} on the left by $\inv(x)$, we obtain
\begin{equation}\label{invxxa_euqals_invxxb}
(\inv(x)\ast x)\ast a = (\inv(x)\ast x)\ast b.
\end{equation}
Since $\inv(x)\ast x = [e,e]$, it follows from \eqref{invxxa_euqals_invxxb} that $a = b$. Hence $\inv(x)\in\W$.}
      
\end{proof}

\end{document}
